\section{Threats to Validity and Reproducibility}
\label{sec:validity}

\subsection{Construct and Internal Validity}
Synthetic scores measure search over fully feasible designed assignments; real gaps depend on user support and utility. Breadth depends on threshold, scale, and grammar, and partial rules may overlap. First attainment describes a trajectory; only the enumerated real spaces provide exact optimality references.

Changing feature selection changes which group predicates can be expressed. Each within-grammar optimizer comparison holds these predicates fixed; the selected-versus-unfiltered ablation intentionally changes them. The original attenuation demonstrates a limitation of the evaluated activity-proxy/backward-elimination wrapper, not an intrinsic failure of all feature selection. Target-by-cap controls test related wrapper configurations; they provide no evidence about active filtering or filter fusion. Minimum support 20 excludes 233 of 294 MovieLens full cells; sparse intersections remain outside that scored audit. Lowering support increases inclusion at the cost of precision. Raising support can remove the highest observed gaps without resolving whether they reflect noise or persistent disadvantage.

The synthetic comparison controls request counts, while the native initializer differs from the challengers' shared initial cells. Relabeling preserves semantic scores but does not couple identical semantic initial candidates across encodings. Encoding spreads therefore combine initialization, search policy, and finite-search variation; operator ablations are still needed for causal attribution to numerical movement. Native source instrumentation and a substituted oracle differ from the original framework's complete execution policy. Different objective geometry and adaptation parameters may change method ordering. Original expanded-audit timings and extension utility-only timings have different scopes and are reported separately.

\subsection{External and Statistical Validity}
The initial experiment covers one landscape; the extension adds three score families with three seeded instances each and repeated encodings/searches. Encodings and search seeds do not increase the number of independent instances. MovieLens separates ten search seeds from three validation-selected NeuMF checkpoints. LFM360K and Amazon each add one supplied DeepFM checkpoint, rather than independent model-training replications. Their 50-candidate utilities, feature schemas, and original preprocessing differ from MovieLens full-catalog ranking, so cross-dataset utility levels are not pooled. Dataset scope was finalized post hoc after checking full-grammar support feasibility; the retained audits do not constitute independent dataset validation. Amazon's history and price profiles do not establish protected-group fairness, and its exact original category/release and history construction are unverified. Demographic labels follow the historical dataset schemas. The 200-attribute work remains a mock integration check. The experiments do not establish general real-model scalability, universal optimizer ordering, or general fairness-signal preservation.

Subgroup rules are adaptively selected using discovery utility. MovieLens feature and checkpoint selection use development users disjoint from the audit population; its test interactions still share audit users and histories with discovery. Additional-dataset confirmation uses disjoint users and fixed discovery rules, but one supplied model and an incompletely reconstructable training protocol remain shared. All extensions are post hoc. Intervals are pointwise, without a familywise multiplicity guarantee; MovieLens bootstrap uncertainty may be understated for small zero-hit cells. The additional-user split does not retroactively preregister hypotheses or establish original model-training separation. Utility gaps alone do not establish causal discrimination, and diagnostic search does not demonstrate mitigation effectiveness.

\subsection{Artifact Reproduction}
The accompanying \path{REPRODUCIBILITY.md} gives exact commands, environments, input requirements, and a directory map separating original training, earlier controls, additional heuristics, and completed additional-dataset audits. Records include source/checkpoint/data hashes, vocabularies, user partitions, exact scores, every requested candidate, and fixed-rule confirmation. Source/configuration checks guard resumption; earlier folders are preserved. CPU runs use an Intel i7-11370H and four inference threads.

The frozen-score exporter verifies model-variable restoration and records library versions, candidate counts, tie handling, and input/output hashes. Modern runners separately execute additional heuristics, external audits, and MovieLens target/cap/support controls. A dedicated revision-asset builder validates request ledgers and regenerates tables without training or search; earlier asset builders alone do not reproduce the revised tables. A dataset-specific manuscript manifest links derived summaries to the preserved source records. All 2.68 million added requests within the manuscript scope were checked for request order, caching, validity, archive counts, and maxima; every real-data score was compared with its exact reference, with deterministic synthetic score spot checks.

The study repository is \url{https://github.com/looppxx/categorical-fairness-search}. The source release includes the search implementations, experiment configurations, reproduction instructions, manuscript sources, and recorded search artifacts. The upstream FairRec source and processed-data entry point is \url{https://github.com/Grey-z/FairRec}. Raw datasets, per-user utility tables, and saved model weights are external inputs whose requirements and provenance are documented in the repository. Hashes identify the original inputs but do not replace access to them. Tables distinguish dataset, checkpoint, search, instance, and encoding repeats.

\section{Conclusion}
\label{sec:conclusion}

This controlled empirical study examines how categorical subgroup search and feature selection affect recommender fairness testing. Its principal contribution is evidence about what an audit can express after projection. In MovieLens, backward elimination from the complete audit schema using activity-proxy labels maps 14, 18, and 15 observed high-gap intersections to projections retaining only 0, 0, and 1 full-cell flags. Target, subset-cap, and support ablations show that predictive relevance and successful optimization are insufficient evidence of intersectional signal retention. These findings concern the evaluated wrapper and do not validate other feature-selection methods.

Five established heuristic challengers broaden the search comparison. GA and ACO exceed native DPSO on the tested synthetic families, while SA has the highest mean maximum on NK/trap instances. In MovieLens, random sampling attains every exact full-grammar optimum and enumeration requires only 294 evaluations; no adaptive-search necessity is demonstrated. Additional LFM360K and Amazon audits show attenuation under a utility-target projection pipeline. Search ordering varies with dataset and grammar, without a universally superior optimizer.

The findings support checking gap retention against an unfiltered reference and evaluating random and exhaustive baselines before attributing benefits to adaptation. Conclusions remain conditional on the evaluated metadata, models, support rules, and ranking protocols. Confirmation intervals include zero, and supplied checkpoint provenance limits the additional-dataset evidence. Richer metadata, independent training replications, comparisons with alternative selection methods, and selection-aware inference are needed for broader effectiveness or population claims. Longitudinal and counterfactual analyses remain separate extensions~\cite{mansoury2020feedback,zoralioglu2024dynamic,li2021towards,medda2024gnnuers}.

